% Principal: GDZ PPN236018647, native 26–28, printed 19–21.
\SourceLeaf{26}\NoSourceHead
\pdfbookmark[1]{Zur Entstehungsgeschichte der Disquisitiones Arithmeticae}{entstehung}
\HistoricTitle{[ZUR ENTSTEHUNGSGESCHICHTE DER DISQUI-\\SITIONES ARITHMETICAE.]}
\begin{center}\large \textsc{Gauss} an v. \textsc{Zimmermann}.\end{center}
\begin{flushleft}\hspace*{.18\linewidth}Hochwohlgeborner Herr\\
\hspace*{.18\linewidth}Verehrenswürdigster Herr Hofrath.\end{flushleft}

Jeder neue Beweis wie sehr Ew. Hochwohlgeb. jede Gelegenheit ergreifen mir nützlich zu sein vermehrt meine Bewunderung und Verehrung Ihrer Güte und bestärkt mich in dem Bestreben mich ihrer würdig zu machen welches ich so viele Ursache habe meine heiligste Pflicht sein zu lassen.

Ich fühle es ganz wie schmeichelhaft mir die Ehre ist wozu Ew. Hochwohlg. mir Hoffnung machen. Gott gebe dem edeln Fürsten ein langes Leben und was können sich nicht die Wissenschaften von ihm versprechen, da er schon eine nur wenige interessirende Arbeit seiner Begünstigung nicht unwerth hält; wie sehr wünschte ich eine gemeinnützigere oder glänzendere geben zu können.

Ich lege hier, Ew. Hochwohlg. Verlangen gemäss einen etwas ausführlichen Plan derselben bei; Sie werden daraus sehen dass der Hauptzweck bloss sein kann den Verstand zu üben und neues Licht über Gegenstände zu verbreiten, die die grössten Geometer unsrer Zeit ihrer eifrigsten Untersuchungen gewürdigt haben. Mir scheint \textsc{Eulers} Ausspruch nicht Unrecht zu haben »Semper cuiusquam problematis, quod a summis ingeniis frustra est tentatum, solutio maximi est momenti«. Und eben dieser grosse Mann hat an mehrern Orten geurtheilt dass Untersuchungen dieser Art zur Übung des Scharfsinns noch dienlicher sind als selbst die Geometrie.
\vfill\begin{flushright}\small 3\textsuperscript{*}\end{flushright}
\clearpage
\SourceLeaf{27}\ArithmeticHead{20}{L}{ARITHMETIK.\quad NACHLASS.}\par\indent Nach meinem Überschlage wird die ganze Schrift gegen 400 Artikel ausmachen wovon im Durchschnitt ungefähr zwei eine gedruckte Quartseite füllen mögte. Daher ich das Ganze etwa auf ein Alphabet schätze; im Octav vielleicht etwas weniger. Dass ich diesen Entwurf nicht zu gross gemacht habe, sondern vielmehr manches dem Leser zu entwickeln zu überlassen genöthigt bin, werden Sie daraus schliessen können, dass ich nothwendig das Wesentliche der Untersuchungen meiner Vorgänger habe mitnehmen müssen, und die von \textsc{Euler} zusammen ungefehr 50, die von \textsc{La Grange} etwa 30 und eine einzige Abhandlung von \textsc{Le Gendre} 12 Bogen beträgt. Ich würde eine Unmöglichkeit unternehmen wenn hierunter nicht manche Wiederholungen wären und durch meine Methoden die weitläuftigsten Rechnungen sehr zusammengezogen würden. Eigentliche Kupfer kann ich entbehren; nur habe ich einen grossen Vorrath von mancherlei Tabellen die ich in müssigen Stunden noch weiter ausdehne und die vielleicht zum Theil sehr nützlich sein können in diesem Felde noch neue Erfindungen zu machen. Von verschiedenen derselben muss ich wenigstens Proben geben. So ist eine Tafel gewiss das bequemste Mittel die Factoren der Zahlen zu bestimmen. Ich füge davon eine kleine Probe bei. Die Punkte bedeuten nichts anders als dass die oben stehende Zahl ein quadratischer Rest der auf der linken Seite stehenden Zahl sei; die Abwesenheit der Punkte das Gegentheil. Ihr Gebrauch würde am bequemsten sein wenn man die einzelnen Streifen A, B, C, D etc. ausschnitte und auf dünne Pappe oder Blech zöge. Es werden dann nach den Umständen gewisse Stäbe zusammengelegt und man sieht gewöhnlich mit der grössten Schnelligkeit welche Zahlen \textls[120]{nicht} Factoren sind. Bei der Schrift selbst wird indessen nur eine kleine Probe der Tafel hinreichen; diess ist nemlich nur Eine einzelne Anwendung der Tafel. Statt der Punkte kann auch jedes andre Zeichen dienen. — — Ich muss noch hinzusetzen dass wenn ich in Absicht des Titels gar nicht genirt bin, ich die Schrift schlechthin Disquisitiones Arithmeticae nennen würde, da ihr Inhalt schwerlich sich genauer würde characterisiren lassen.

Der Hr. Hofrath \textsc{Kästner} hat mir versprochen mir die Bekanntschaft des Hn. P. \textsc{Hindenburg} zu verschaffen; ich habe ihm einen Aufsatz zur Ansicht vorgelegt, und ich denke denselben wenn er mir sein Urtheil darüber wird gesagt haben an \textsc{Hindenburg} zu schicken.
\clearpage
\SourceLeaf{28}\ArithmeticHead{21}{R}{GAUSS AN V. ZIMMERMANN.}\par\indent Es ist mir ein sehr angenehmes Geschäft Ew. Hochwohlgeb. von meinem jungen Schweizer Nachricht zu geben. Er ist aus Thun im Canton Bern und hat sich schon in Bern unter \textsc{Tralles} Anleitung gute Mathematische Kenntnisse erworben, die er hier seit einem halben Jahre mit dem besten Erfolg erweitert. Er wird hier bis Michaelis bleiben und dann nach Paris gehen, ehe er Deutschland verlässt denkt er noch mehrere Örter davon zu besuchen, und er wird es nicht versäumen auch Braunschweig zu sehen. Sein Name ist \textsc{Beckh}.

Ich bin glücklich mich mit der innigsten Verehrung und Dankbarkeit nennen zu können
\begin{flushright}
Hochwohlgeborner Herr Hofrath\\[.4em]
Ew. Hochwohlgeborn\\[.4em]
ergebensten Diener
\end{flushright}
\noindent Göttingen den 12\textsuperscript{ten} Merz 1797.\hfill\textsc{Gauss}.

\bigskip\begin{center}\small\bfseries BEMERKUNG.\end{center}
\begingroup\small\par\indent
Der hier mit der ursprünglichen Rechtschreibung abgedruckte Brief befindet sich gegenwärtig im Besitz der Königlichen Bibliothek zu Berlin. Die Aufschrift lautet
\begin{center}
\begin{tabular}{ll}
\multicolumn{2}{l}{Se. Hochwohlgeborn}\\
&dem Herrn Hofrath\\
&\quad v. \textsc{Zimmermann}\\
durch Einschluss.&\qquad Braunschweig
\end{tabular}
\end{center}
\noindent Zwei spätere Briefe von \textsc{Gauss} an v. \textsc{Zimmermann} vom 22. November und vom 24. Dezember 1797, die sich ebenfalls auf die \textit{Disqu. Arithm.} beziehen, sind im Auszug veröffentlicht in der Schrift von L. \textsc{Hänselmann}, \textit{Karl Friedrich Gauss. Zwölf Kapitel aus seinem Leben}, Leipzig 1878, S. 34--37. Vergl. den Artikel 2 des \textsc{Bachmann}schen Aufsatzes »Über \textsc{Gauss}' zahlentheoretische Arbeiten«.

Den S. 19, Zeile 4, 3 v. u. angeführten Ausspruch \textsc{Eulers} hat \textsc{Gauss} auf das vordere Schutzblatt seines Exemplars von \textsc{Lamberts} \textit{Zusätzen} eingetragen, mit der Quellenangabe: \textsc{Euler}, Comm. N. Petr. XI, p. 153; er findet sich in der Tat in der Abhandlung: \textit{De motu corporis ad duo centra virium fixa attracti}, Novi comm. Acad. Petrop. \textbf{11} (1765) 1767, S. 152--184, auf S. 153. Das S. 19, Z. 2 v. u. erwähnte Urteil \textsc{Eulers} lautet: »Quare cum vulgo ad ratiocinii facultatem comparandam demonstrationes geometricae commendari soleant, quippe quae regularum ratiocinandi usum maxime contineant, nescio an non ad hunc scopum demonstrationes arithmeticae multo magis sint accommodatae: in his enim multo maiori cura est cavendum, ne a praescriptis Logicorum regulis aberremus, quoniam plerumque nimis est difficile, in errorem non prolabi« und steht in der Abhandlung: \textit{Specimen de usu observationum in mathesi pura}, Novi comment. Acad. Petrop. \textbf{6} (1756/7), 1761, S. 185--230, auf S. 187. Die S. 20, Zeile 18 erwähnte »kleine Probe« der zur Bestimmung der Faktoren von Zahlen dienenden Tafel findet sich in den \textit{Disqu. Arithm.}, Werke I, S. 469, die Anleitung zu ihrem Gebrauch ebenda, S. 404, 405.
\begin{flushright}\textsc{Schlesinger}.\end{flushright}
\endgroup
\ClosingRule
\clearpage
